% Generated draft: requires mathematical and source review.
\documentclass[11pt]{article}
\usepackage[margin=25mm]{geometry}
\usepackage{fontspec}
\setmainfont{Latin Modern Roman}
\usepackage{amsmath,amssymb,mathtools,unicode-math}
\setmathfont{Latin Modern Math}
\usepackage{enumitem,needspace,xurl,hyperref}
\setlength{\parindent}{0pt}
\setlength{\parskip}{5pt}
\setlength{\emergencystretch}{3em}
\newcommand{\N}{\mathbb{N}}
\newcommand{\Z}{\mathbb{Z}}
\newcommand{\Q}{\mathbb{Q}}
\newcommand{\R}{\mathbb{R}}
\newcommand{\C}{\mathbb{C}}
\newcommand{\F}{\mathbb{F}}
\newcommand{\A}{\mathbb{A}}
\newcommand{\PP}{\mathbb{P}}
\newcommand{\E}{\mathbb{E}}
\newcommand{\eps}{\varepsilon}
\newcommand{\dd}{\,\mathrm{d}}
\begin{document}
\section*{UnsolvedMath ID: 10400035}
\textbf{Source problem code:} catalogue entry
\subsection{Definitions and mathematical statement}
An \(m\)-strand string link is a collection of disjoint, properly embedded oriented arcs in a \(3\)-ball, with prescribed endpoints and strands indexed from bottom to top; equivalence and boundary fixing are understood in the source's category. A rational finite-type invariant of order at most \(r\) extends by the crossing-difference rule to singular string links and vanishes on those with more than \(r\) transverse double points. The record does not specify whether its \(V_n\) means this filtered space, an exact-degree graded quotient, or another convention, nor whether the number of strands is fixed. Retain \(V_n\) with that unresolved source meaning. Let \(M_n\subset V_n\) be the source's Milnor-invariant subspace; its precise inclusion of products and degree bounds must be checked.

T. Stanford's Question 2.13 asks whether \(M_n\) has an interesting complementary space in \(V_n\). Specifically, define
\[
 N_n=\{v\in V_n:\ v(L)=0\text{ for every string link }L
       \text{ with }\pi_1(B^3\setminus L)\text{ free}\}.
\]
Is \(N_n\) nonzero? Do \(N_n\) and \(M_n\) span \(V_n\), that is,
\[
 V_n=M_n+N_n\ ?
\]
The record gives no range of \(n\) or component counts for these questions. The displayed spanning question does not additionally assert \(M_n\cap N_n=\{0\}\); a direct-sum complement would require that further condition.

The supplied motivation reports that Milnor invariants are well defined on string links, whereas for links their usual definition has indeterminacy; they are finite-type invariants. It reports that Milnor invariants, including products, correspond on the Jacobi-diagram level to functionals vanishing on diagrams with internal loops, and are the rational-valued finite-type concordance invariants of string links. It also states that local knots in strands are not detected by Milnor invariants and that their presence makes the complement group nonfree, motivating detection of nonfreeness by finite-type invariants. These contextual claims retain the source's scope and require checking its conventions.

M. Polyak's supplied contextual description associates \(\mu_{i_1\cdots i_k,n}(L)\) to the coefficient of \(X_{i_1}\cdots X_{i_k}\) in a Magnus expansion with noncommuting variables, formally substituting \(m_j=1+X_j\) in a meridian expression for the last component. It requires vanishing of lower Milnor invariants for the stated link invariance. The examples given are \(\mu_{1,2}=\operatorname{lk}(L_1,L_2)\), and, under vanishing pairwise linking numbers, \(\mu_{12,3}=\operatorname{lk}(L_{12},L_3)\). More generally the excerpt writes
\[
 \mu_{12\cdots n-1,n}=\operatorname{lk}(L_{12\cdots n-1},L_n),
\]
under vanishing of lower invariants, with \(L_{12\cdots k}\) obtained recursively as intersections of Seifert surfaces for \(L_{12\cdots k-1}\) and \(L_k\), \(2\le k\le n-1\). This compressed account omits meridian-word/nilpotent-quotient, orientation, and surface-choice details; it is preserved as motivation, not adopted as a complete construction valid for arbitrary surface choices.
\subsection{Short English statement}
Among finite-type string-link invariants, ask whether those vanishing whenever the complement group is free form a nonzero space and, together with the Milnor invariants, span the full space in question. The source motivates this by local knotting invisible to Milnor invariants. The degree convention, strand count, and intended meaning of complement remain to be fixed.
\subsection{Sources}
\begin{itemize}
\item Local UnsolvedMath registry, permanent ID \(10400035\), exact source code \texttt{catalogue entry}; title truncated to ``Question 2.13 --- (T.'' The statement attributes the main question to T. Stanford and the subsequent Milnor-invariant explanation to M. Polyak.
\item Ohtsuki, \emph{Problems on invariants of knots and 3-manifolds} (2002), Question 2.13, PDF page 36: \url{https://msp.org/gtm/2002/04/gtm-2002-04-024s.pdf}. Bold headings distinguish 231 numbered problems, conjectures, and questions from surrounding references.
\item Supplied internal references: Milnor [283] (1954); Habegger and Lin [163]; Bar-Natan [26] and Lin [261] independently on finite type; Habegger and Masbaum [164]; and [89] for the reviewed Milnor construction. Complete bibliography entries are absent and have not been invented.
\item Extraction pdf-font-structure, accessed 2026-07-29. Source evidence NEEDS\_REVIEW presents an open problem and requires a current primary-literature audit. The supplied OPEN-TRIAGE summary calls the complementary-space question not obviously resolved.
\item Dataset attribution: Ulam AI / UnsolvedMath Contributors, CC BY 4.0. Source rights NEEDS\_REVIEW: publicly accessible author/publisher copy, redistribution terms requiring review.
\end{itemize}
\end{document}
